El presente proyecto se enfocó en la etapa del proceso comercial donde los leads son activamente gestionados y priorizados, utilizando la etapa de reserva como punto de corte para evaluar los resultados de conversión. Dentro de este límite, el trabajo se centró en establecer criterios estructurados y basados en datos que sirvan de apoyo a la manera en que los leads son evaluados y priorizados en proyectos residenciales de vivienda de empresas desarrolladoras en la Ciudad de Guatemala. De esta manera, el alcance del estudio se delimita a una fase específica del embudo comercial, sin extenderse a etapas posteriores del proceso de venta.

Desde una perspectiva de viabilidad, el proyecto dependió del acceso a datos comerciales existentes y del conocimiento de los procesos internos vinculados a proyectos residenciales específicos. Los recursos requeridos fueron primordialmente digitales, incluyendo conjuntos de datos derivados de sistemas CRM. Estos recursos estuvieron disponibles a través de la coordinación con la entidad que facilitó la información, lo que permitió realizar el análisis sin necesidad de infraestructura adicional ni de sistemas nuevos de recolección de datos. Las fuentes principales incluyeron datos internos, documentación del proceso comercial y hallazgos recuperados de las partes interesadas involucradas en la gestión de leads, lo que permitió una comprensión fundamentada de cómo se toman actualmente las decisiones comerciales. Los recursos secundarios consistieron en literatura académica, reportes de la industria e investigación previa en calificación de leads, modelado predictivo y optimización del proceso comercial, que proveyeron la base conceptual y metodológica para el análisis.

En términos de tiempo, el proyecto se estructuró para alinearse con el calendario académico, asignando tiempo a cada una de las fases, desde el entendimiento del negocio hasta el desarrollo del modelo y la propuesta de implementación. Esta planificación permitió que cada etapa contara con espacio suficiente para su desarrollo y revisión, y que el avance del trabajo se monitoreara de manera ordenada y consistente.

El estudio está sujeto a diversas limitaciones, varias de las cuales se materializaron durante el desarrollo y condicionaron los resultados obtenidos.

En primer lugar, la disponibilidad y calidad de los datos alteró el alcance del análisis. La ausencia de información resultó ser masiva y estructurada: las variables de perfil demográfico presentan porcentajes de valores faltantes que alcanzan el 93.5\%, dado que no se registran en el primer contacto sino conforme avanza la calificación del lead. Esta característica obligó a redefinir el universo de modelado y a tratar la ausencia como categoría informativa en lugar de imputarla. Adicionalmente, se identificaron trece columnas con fuga de información, campos que solo se completan una vez cerrada la venta, cuya remoción redujo sustancialmente el desempeño aparente de las primeras versiones del modelo.

En segundo lugar, el proyecto no se extendió más allá de la etapa de reserva, excluyendo procesos posteriores tales como la formalización del contrato, la promesa de compra o el análisis de crédito.

En tercer lugar, aunque el estudio se plantea desde una perspectiva metodológica aplicable al contexto de las empresas desarrolladoras inmobiliarias en la Ciudad de Guatemala, el análisis empírico se desarrolló con información correspondiente a seis proyectos residenciales específicos, por lo que los hallazgos deben interpretarse dentro de ese alcance. Para mitigar parcialmente esta limitación, la identidad del proyecto se excluyó deliberadamente del conjunto de predictores, con un costo estadístico cuantificado de aproximadamente 5\% en el desempeño, a cambio de que el modelo sea aplicable a proyectos residenciales nuevos.

En cuarto lugar, los modelos se desarrollaron en un entorno analítico controlado y no fueron implementados en sistemas en producción, lo cual restringe la capacidad de medir el impacto en tiempo real. Todos los resultados provienen de validación cruzada sobre datos históricos: no se dispuso de una ventana temporal posterior ni de una prueba controlada que permitiera comparar el resultado comercial con y sin priorización. La ganancia reportada es, por lo tanto, una estimación retrospectiva de la calidad del ordenamiento y no una medición de impacto en ventas.

Finalmente, el modelo reordena la cartera comercial pero no identifica compradores con certeza. Esta precisión delimita el tipo de decisión que el modelo puede sostener.

En cuanto al uso de los datos, el proyecto se condujo bajo autorización para trabajar con información real con propósitos académicos. Cualquier dato sensible se manejó bajo acuerdos de confidencialidad y los resultados se presentan de manera que no expongan información personal identificable ni detalles comerciales sensibles. Adicionalmente, aunque las variables de sexo y estado civil se emplearon con finalidad descriptiva en el perfilamiento de segmentos, se establece explícitamente que no deben utilizarse como criterio de asignación de esfuerzo comercial.